\section{Limits of Quantum Information}

\subsection{No-Cloning Principle}

\subsubsection{Why Cloning Matters}

Briefly, the \important{no-cloning principle} states that it is \textbf{impossible to create a perfect copy of an arbitrary unknown quantum state}.

\highspace
At first sight this may seem like a strange technical limitation, because in classical computing copying information is \emph{trivial}: we can easily copy files, duplicate RAM, clone hard disks, send packets over the network, etc. However, \hl{in quantum computing, copying an arbitrary qubit is fundamentally impossible.} This is \textbf{not a technological limitation}, instead it is a direct \textbf{consequence of the mathematical structure} of quantum mechanics.

\highspace
In classical computing, a bit can be in one of two states: \texttt{0} or \texttt{1}. In quantum computing, a qubit instead can be in infinitely many states, because it can be in any superposition of \texttt{0} and \texttt{1}:
\begin{equation*}
    \ket{\psi} = a\ket{0} + b\ket{1}
\end{equation*}
With complex amplitudes $a$ and $b$ such that $|a|^2 + |b|^2 = 1$. So a qubit is not just ``0'' or ``1'', but it represents a point on the Bloch sphere, meaning there are infinitely many possible states. \hl{If arbitrary cloning were possible}, we could take an unknown state $\ket{\psi}$ and create as many copies of it as we want. This would allow us to perform measurements on the copies to learn about the original state $\ket{\psi}$ without disturbing it, \textbf{violating the fundamental principles of quantum mechanics}.

\highspace
\begin{flushleft}
    \textcolor{Red2}{\faIcon{link} \textbf{Relation with Unitary}}
\end{flushleft}
The no-cloning principle theorem is deeply connected to the fact that quantum evolution is \textbf{unitary}. Quantum gates are represented by unitary matrices:
\begin{equation*}
    U^\dagger U = U U^\dagger = I
\end{equation*}
Unitary operations preserve norms, probabilities and inner products between states. \textbf{Cloning would violate this preservation property}. Because, intuitively, if two partially similar states must remain partially similar after evolution (i.e., their inner product must be preserved), then cloning would instead make them ``more distinguishable'' (i.e., their inner product would change), which is a contradiction. We will see the formal proof of the no-cloning principle in the next section.

\highspace
\begin{flushleft}
    \textcolor{Red2}{\faIcon{link} \textbf{Relation with No-Signaling}}
\end{flushleft}
The no-cloning principle is also related to the no-signaling principle, because \textbf{breaking the no-cloning principle would violate the no-signaling principle}.

\highspace
Suppose Alice and Bob share an entangled Bell pair:
\begin{equation*}
    \ket{\Phi^+} = \frac{1}{\sqrt{2}}(\ket{00} + \ket{11}) = \frac{\ket{00} + \ket{11}}{\sqrt{2}}
\end{equation*}
Normally, Bob cannot determine whether Alice has measured her qubit or not, because his local measurements always appear random (50\% probability of measuring \texttt{0} and 50\% probability of measuring \texttt{1}). But \emph{imagine} Bob could clone his qubit infinitely many times. Then, he could:
\begin{enumerate}
    \item Create many identical copies of his qubit.
    \item Perform measurements/statistics on all copies.
    \item Infer whether the global quantum state collapsed (i.e., whether Alice measured her qubit). Using the principle of reductio ad absurdum, if Alice does not measure, then Bob will measure random results (0, 1, 0, 1, etc.). When Alice measures and collapses the state to \ket{0}, then Bob's qubit will instantly become \ket{0}, and therefore, all of his copies will be \ket{0}.
\end{enumerate}
This would allow instantaneous communication between Alice and Bob, faster than the speed of light, which is forbidden by the no-signaling principle (violation of relativity, no information can travel faster than the speed of light). Therefore, the \textbf{impossibility of cloning helps preserve the no-signaling principle}.

\highspace
\begin{flushleft}
    \textcolor{Red2}{\faIcon{link} \textbf{Relation with Quantum State Tomography}}
\end{flushleft}
Another consequence of breaking the no-cloning principle would be that we would be \hl{able to store an infinite amount of classical information inside one qubit.}

\highspace
The idea is the following. Because a qubit can assume infinitely many states, one could theoretically encode a very long classical string into amplitudes (e.g., the binary expansion of $\pi$):
\begin{equation*}
    \ket{\psi_{\text{message}}}
\end{equation*}
But there is a problem: we \textbf{cannot directly read the amplitudes} from a single qubit, because measurement collapses the state. If we measure the qubit, we only get one classical outcome (e.g., \texttt{0} or \texttt{1}), and not the full amplitudes $a$ and $b$. So normally, the information is not accessible.

\highspace
However, exists a technique called \important{quantum state tomography} that:
\begin{enumerate}
    \item Prepares many identical copies\footnote{
        Quantum State Tomography does not violate the no-cloning principle, because it does not create perfect copies of an unknown state. Instead, it prepares many identical copies of a \textbf{known} state $\ket{\psi_{\text{message}}}$, which is allowed by quantum mechanics. For example, if we have a quantum circuit that prepares the state $\ket{\psi_{\text{message}}}$, we can run that circuit many times to get many copies of $\ket{\psi_{\text{message}}}$, which is perfectly fine. The no-cloning principle only forbids creating perfect copies of an \textbf{unknown} state, not preparing many copies of a \textbf{known} state.
    } of the same state $\ket{\psi_{\text{message}}}$.
    \item Measures them in different bases (e.g., X, Y, Z).
    \item Reconstructs the amplitudes of the original state $\ket{\psi_{\text{message}}}$ statistically.
\end{enumerate}
In simple terms, given many copies of the same unknown state, estimate what the state was.

\highspace
Now, imagine there were \textbf{no no-cloning theorem}. Then we could do the following:
\begin{enumerate}
    \item Take the single qubit $\ket{\psi_{\text{message}}}$ that encodes the message.
    \item Clone it infinitely many times to get many copies of $\ket{\psi_{\text{message}}}$. This would be possible because we are assuming that the no-cloning principle does not hold.
    \item Perform measurements/statistics on all copies to reconstruct the amplitudes of the original state $\ket{\psi_{\text{message}}}$, and therefore read the message encoded in the amplitudes.
\end{enumerate}
So cloning would transform ONE unknown qubit into enough data for tomography, which would allow us to read the message encoded in the amplitudes. This would be a huge problem, because we could encode an infinite amount of classical information into one qubit and then read it using cloning and tomography.

\begin{examplebox}[: Why Cloning Matters]
    \textcolor{Green3}{\faIcon{check-circle} \textbf{Normal world (no-cloning theorem exists).}} Alice wants to send Bob one qubit:
    \begin{equation*}
        \ket{\psi_{\text{message}}} = \sqrt{0.73148291}\ket{0} + \sqrt{0.26851709}\ket{1}
    \end{equation*}
    Suppose these decimals secretly encode a super secret message. Bob receives the qubit. Now Bob asks: ``\emph{what are the exact amplitudes of the state $\ket{\psi_{\text{message}}}$?}'' Bob measures the qubit, but measurement only gives him one classical bit (0 or 1). For example, if Bob measures $0$, then the state collapses to $\ket{0}$, and the original amplitudes are lost. So from ONE measurement, Bob cannot reconstruct $0.73148291$ exactly. So the hidden information is inaccessible to Bob, and the message is safe.

    \highspace
    \textcolor{Red2}{\faIcon{exclamation-triangle} \textbf{Hypothetical world (no-cloning theorem does not exist).}} Now suppose cloning is allowed. Bob receives the qubit $\ket{\psi_{\text{message}}}$. Now Bob \textbf{can clone} the qubit, so he runs a mysterious cloning machine that creates many identical clones of $\ket{\psi_{\text{message}}}$:
    \begin{equation*}
        \ket{\psi} \rightarrow \ket{\psi}\ket{\psi}\ket{\psi}\ket{\psi}\ket{\psi} \cdots
    \end{equation*}
    Now he has millions of copies of $\ket{\psi_{\text{message}}}$. Now Bob can decrypt the super secret message. How? Simple:
    \begin{enumerate}
        \item \textbf{Measure many copies in the computational basis}. Bob has millions of copies. Suppose he observes (after many measurements) that 731482 times he gets $0$ and 268517 times he gets $1$. So Bob can estimate the probabilities of measuring $0$ and $1$ as $0.731482$ and $0.268517$. Bob is already very close to the original amplitudes, but he can do better.
        \item \textbf{Measure in other bases}. Bob is smart, and he knows that measuring in the computational basis only gives him the magnitudes of the amplitudes, but not their phases. So Bob also measures in the X and Y bases to get more information about the state. By combining all the measurement results, Bob can reconstruct the original state $\ket{\psi_{\text{message}}}$ with very high precision, and therefore read the super secret message encoded in the amplitudes.
    \end{enumerate}
    \textcolor{Red2}{\faIcon{exclamation-triangle} \textbf{The Problem.}} Alice effectively transmitted thousands, millions, theoretically infinite classical digits using only one qubit, that should not be physically possible (because in physics, the amount of extractable classical information from a physical system must be finite).
\end{examplebox}

\begin{flushleft}
    \textcolor{Red2}{\faIcon{link} \textbf{Relation with Quantum Error Correction}}
\end{flushleft}
In classical systems, redundancy is simple: we can copy bits to create redundancy for error correction. For example, we can encode a bit as three bits:
\begin{equation*}
    0 \rightarrow 000 \quad\quad 1 \rightarrow 111
\end{equation*}
Then, if one bit gets flipped, we can still recover the original bit by majority voting:
\begin{equation*}
    000 \rightarrow 000 \quad\quad 001 \rightarrow 000 \quad\quad 010 \rightarrow 000 \quad\quad 100 \rightarrow 000
\end{equation*}
Then majority voting allows us to correct single bit errors. However, in \hl{quantum computing we cannot simply duplicate qubits} and we cannot create copies of arbitrary quantum states. So \textbf{quantum error correction cannot protect information by simply copying qubits}. Instead, it uses more complex techniques (e.g., entanglement, syndrome measurements, etc.) to protect quantum information without violating the no-cloning principle.

\highspace
In summary, the no-cloning theorem is not an isolated curiosity. It is one of the foundational constraints that shape all of quantum information theory. It is connected to reversibility, unitarity, causality, measurement, entanglement, quantum cryptography, quantum error correction. \hl{Without the no-cloning theorem, the entire structure of quantum information would collapse.}